\section{AMC 12 12A 2007}
        \begin{enumerate}
        	\item (\textit{AMC 12 12A 2007}) One ticket to a show costs $\$20$ at full price. Susan buys $4$ tickets using a coupon that gives her a $25\%$ discount. Pam buys $5$ tickets using a coupon that gives her a $30\%$ discount. How many more dollars does Pam pay than Susan?


 $\textbf{(A)} \ 2 \qquad \textbf{(B)} \ 5 \qquad \textbf{(C)} \ 10 \qquad \textbf{(D)} \ 15 \qquad \textbf{(E)} \ 20$

 \textbf{Answer: }C



\textbf{Solution}

$\textbf{Answer: \boxed{(C)}}$ Susan pays $(4)(0.75)(20) = 60$ dollars. Pam pays $(5)(0.70)(20) = 70$ dollars, so she pays $70-60=10$ more dollars than Susan.

 


	\item (\textit{AMC 12 12A 2007}) An aquarium has a \href{https://artofproblemsolving.com/wiki/index.php?title=Rectangular\_prism}{rectangular base} that measures 100 cm by 40 cm and has a height of 50 cm. It is filled with water to a height of 40 cm. A brick with a rectangular base that measures 40 cm by 20 cm and a height of 10 cm is placed in the aquarium. By how many centimeters does the water rise?


 $\mathrm{(A)}\ 0.5\qquad \mathrm{(B)}\ 1\qquad \mathrm{(C)}\ 1.5\qquad \mathrm{(D)}\ 2\qquad \mathrm{(E)}\ 2.5$

 \textbf{Answer: }D



\textbf{Solution}

The brick has volume $8000 cm^3$. The base of the aquarium has area $4000 cm^2$. For every cm the water rises, the volume increases by $4000 cm^3$; therefore, when the volume increases by $8000 cm^3$, the water level rises $2 cm \Rightarrow\fbox{D}$



	\item (\textit{AMC 12 12A 2007}) The larger of two consecutive odd integers is three times the smaller. What is their sum?


 $\mathrm{(A)}\ 4\qquad \mathrm{(B)}\ 8\qquad \mathrm{(C)}\ 12\qquad \mathrm{(D)}\ 16\qquad \mathrm{(E)}\ 20$

 \textbf{Answer: }A



\textbf{Solution 1}

Let $n$ be the smaller term. Then $n+2=3n \Longrightarrow 2n = 2 \Longrightarrow n=1$

            \begin{itemize}
        \item Thus, the answer is $1+(1+2)=4 \mathrm{(A)}$
            \end{itemize}
            




\textbf{Solution 2}


            \begin{itemize}
        \item By trial and error, 1 and 3 work. 1+3=4.
            \end{itemize}
            


	\item (\textit{AMC 12 12A 2007}) Kate rode her bicycle for 30 minutes at a speed of 16 mph, then walked for 90 minutes at a speed of 4 mph. What was her overall average speed in miles per hour?


 $\mathrm{(A)}\ 7\qquad \mathrm{(B)}\ 9\qquad \mathrm{(C)}\ 10\qquad \mathrm{(D)}\ 12\qquad \mathrm{(E)}\ 14$

 \textbf{Answer: }A



\textbf{Solution}


 \[16 \cdot \frac{30}{60}+4\cdot\frac{90}{60}=14\]
 \[\frac{14}2=7\Rightarrow\boxed{A}\]



	\item (\textit{AMC 12 12A 2007}) Last year Mr. Jon Q. Public received an inheritance. He paid $20\%$ in federal taxes on the inheritance, and paid $10\%$ of what he had left in state taxes. He paid a total of $\textdollar10,500$ for both taxes. How many dollars was his inheritance?


 $(\mathrm {A})\ 30\,000 \qquad (\mathrm {B})\ 32\,500 \qquad (\mathrm {C})\ 35\,000 \qquad (\mathrm {D})\ 37\,500 \qquad (\mathrm {E})\ 40\,000$

 \textbf{Answer: }D



\textbf{Solution 1}

After paying his taxes, he has $0.8*0.9=0.72$ of his earnings left. Since $10500$ is $0.28$ of his income, he got a total of $\frac{10500}{0.28}=37500\  \mathrm{(D)}$.

 



\textbf{Solution 2}

Let his total inheritance be a number $x$. The amount of tax he has to pay at first can be represented as $0.2x$. The amount of money he has left over after this expense can be written as $0.8x$. He then has to pay $10\%$ tax on this remaining amount, which can be written as $0.1(0.8x) = 0.08x$. The combined expense of these taxes is $10500$, so $0.2x + 0.08x = 10500 \Longrightarrow 0.28x = 10500$. Therefore $x = \frac{10500}{0.28}\ = 37500  \mathrm{(D)}$. \href{https://artofproblemsolving.com/wiki/index.php?title=User:Ankitamc&action=edit&redlink=1}{Ankitamc} (\href{https://artofproblemsolving.com/wiki/index.php?title=User\_talk:Ankitamc&action=edit&redlink=1}{talk}) 13:32, 17 January 2021 (EST)AnkitAmc.~edited by mobius247

 


	\item (\textit{AMC 12 12A 2007}) Triangles $ABC$ and $ADC$ are \href{https://artofproblemsolving.com/wiki/index.php?title=Isosceles}{isosceles} with $AB=BC$ and $AD=DC$. Point $D$ is inside triangle $ABC$, angle $ABC$ measures 40 degrees, and angle $ADC$ measures 140 degrees. What is the degree measure of angle $BAD$?


 $\mathrm{(A)}\ 20\qquad \mathrm{(B)}\ 30\qquad \mathrm{(C)}\ 40\qquad \mathrm{(D)}\ 50\qquad \mathrm{(E)}\ 60$

 \textbf{Answer: }D



\textbf{Solution 1}


                    \begin{figure}[H]
                        \centering
                        \includegraphics[width=0.5\textwidth]{Solution/AMC_12_12A_2007_6_Solution_Figure_1.png}
                    \end{figure}
                    

 We angle chase and find out that:

 
            \begin{itemize}
        \item $\angle DAC=\frac{180-140}{2} = 20$
        \item $\angle BAC=\frac{180-40}{2} = 70$
        \item $\angle BAD=\angle BAC- \angle DAC=50\ \mathrm{(D)}$
            \end{itemize}
            



\textbf{Solution 2}


                    \begin{figure}[H]
                        \centering
                        \includegraphics[width=0.5\textwidth]{Solution/AMC_12_12A_2007_6_Solution_Figure_2.png}
                    \end{figure}
                    

 Since triangle $ABC$ is isosceles we know that angle $\angle BAC = \angle BCA$.

 Also since triangle $ADC$ is isosceles we know that $\angle DAC = \angle DCA$.

 This implies that  $\angle BAD = \angle BCD$.

 Then the sum of the interior angles of quadrilateral $ABCD$ is  $40 + 220 + 2\angle BAD = 360$.

 Solving the equation, we get $\angle BAD = 50$.

 Therefore the answer is $\mathrm{(D)}$.

 



\textbf{Solution 3}

Using the previous drawing we know that ACD and DAC are both equal to 20 (40/2 or (180-40)/2). We also know that BAC and BCA are both 70 or (180-40)/2. Thus BDA is 70-20 which is equal to 50. 

 -RealityWrites

 


	\item (\textit{AMC 12 12A 2007}) Let $a, b, c, d$, and $e$ be five consecutive terms in an arithmetic sequence, and suppose that $a+b+c+d+e=30$. Which of $a, b, c, d,$ or $e$ can be found?


 $\textrm{(A)} \ a\qquad \textrm{(B)}\ b\qquad \textrm{(C)}\ c\qquad \textrm{(D)}\ d\qquad \textrm{(E)}\ e$

 \textbf{Answer: }C



\textbf{Solution}

Let $f$ be the common difference between the terms.

 
            \begin{itemize}
        \item $a=c-2f$
        \item $b=c-f$
        \item $c=c$
        \item $d=c+f$
        \item $e=c+2f$
            \end{itemize}
            $a+b+c+d+e=5c=30$, so $c=6$. But we can't find any more variables, because we don't know what $f$ is. So the answer is $\textrm{C}$.

 


	\item (\textit{AMC 12 12A 2007}) A star-\href{https://artofproblemsolving.com/wiki/index.php?title=Polygon}{polygon} is drawn on a clock face by drawing a \href{https://artofproblemsolving.com/wiki/index.php?title=Chord}{chord} from each number to the fifth number counted clockwise from that number. That is, chords are drawn from 12 to 5, from 5 to 10, from 10 to 3, and so on, ending back at 12. What is the degree measure of the \href{https://artofproblemsolving.com/wiki/index.php?title=Angle}{angle} at each \href{https://artofproblemsolving.com/wiki/index.php?title=Vertex}{vertex} in the star polygon?


 $\mathrm{(A)}\ 20\qquad \mathrm{(B)}\ 24\qquad \mathrm{(C)}\ 30\qquad \mathrm{(D)}\ 36\qquad \mathrm{(E)}\ 60$

 \textbf{Answer: }C



\textbf{Solution}


                    \begin{figure}[H]
                        \centering
                        \includegraphics[width=0.5\textwidth]{Solution/AMC_12_12A_2007_8_Solution_Figure_1.png}
                    \end{figure}
                    

 We look at the angle between 12, 5, and 10. It subtends $\frac 16$ of the circle, or $60$ degrees (or you can see that the \href{https://artofproblemsolving.com/wiki/index.php?title=Arc}{arc} is $\frac 23$ of the \href{https://artofproblemsolving.com/wiki/index.php?title=Right\_angle}{right angle}). Thus, the angle at each vertex is an \href{https://artofproblemsolving.com/wiki/index.php?title=Inscribed\_angle}{inscribed angle} subtending $60$ degrees, making the answer $\frac 1260 = 30^{\circ} \Longrightarrow \mathrm{(C)}$



	\item (\textit{AMC 12 12A 2007}) Yan is somewhere between his home and the stadium. To get to the stadium he can walk directly to the stadium, or else he can walk home and then ride his bicycle to the stadium. He rides 7 times as fast as he walks, and both choices require the same amount of time. What is the \href{https://artofproblemsolving.com/wiki/index.php?title=Ratio}{ratio} of Yan's distance from his home to his distance from the stadium?


 $\mathrm{(A)}\ \frac 23\qquad \mathrm{(B)}\ \frac 34\qquad \mathrm{(C)}\ \frac 45\qquad \mathrm{(D)}\ \frac 56\qquad \mathrm{(E)}\ \frac 78$

 \textbf{Answer: }B



\textbf{Solution 1}

Let the distance from Yan's initial position to the stadium be $a$ and the distance from Yan's initial position to home be $b$. We are trying to find $b/a$, and we have the following identity given by the problem:

 
 \begin{align*}a &= b + \frac{a+b}{7}\\ \frac{6a}{7} &= \frac{8b}{7} \Longrightarrow 6a = 8b\end{align*}


 Thus $b/a = 6/8 = 3/4$ and the answer is $\mathrm{(B)}\ \frac{3}{4}$

 Note: The identity can be given from an application of $t=\frac dr$.

 



\textbf{Solution 2}

Another way of solving this problem is by setting the distance between Yan's home and the stadium, thus filling in one variable.Let us set the distance between the two places to be $42z$, where $z$ is a random measurement (cause life, why not?). The time going to his home then riding his bike, which is $7$ times faster, to the stadium is equal to the time it takes for him to walk to the stadium. So the equation would be:Let $x=$ the distance from Yan's position to his home. Let $42z=$ the distance from Yan's home to the stadium.

 $42z-x=x+\frac{42z}{7}$
$42z-x=x+6z$
$36z=2x$
$x=18z$

 But we're still not done with the question. We know that Yan is $18z$ from his home, and is also $42z-18z$ or $24z$ from the stadium.$\frac{18z}{24z}$  , the $z$'s cancel out, and we are left with $\frac{3}{4}$.Thus, the answer is $\mathrm{(B)}\ \frac{3}{4}$




\textbf{Solution 3}

Assume that the distance from home to the stadium is 1, and the distance from Yan to home is $b$. Also assume that the speed of walking is 1, so the speed of biking is 7. Thus$1-b=b+\frac{1}{7} \Longrightarrow b=\frac{3}{7}$.We want to find $\frac{b}{1-b}$ which is equal to

 $\frac{3}{7} \div \frac{4}{7} = \mathrm{(B)}\ \frac{3}{4}$




\textbf{Solution 4}


                    \begin{figure}[H]
                        \centering
                        \includegraphics[width=0.5\textwidth]{Solution/AMC_12_12A_2007_9_Solution_Figure_1.png}
                    \end{figure}
                    Let $H$ represent Yan's home, $S$ represent the stadium, and $Y$ represent Yan's current position. If Yan walks directly to the stadium, he will reach Point $P$ the same time he will reach $H$ if he is walking home. Since he bikes $7$ times as fast as he walks and the time is the same, the distance from his home to the stadium must be $7$ times the distance from $P$ to the stadium. If $PS=x$, then $HS=7x$ and $HP=6x$. Since Y is the midpoint of $\overline{HP}$, $HY=YP=3x$. Therefore, the ratio of Yan's distance from his home to his distance from the stadium is $\frac{YH}{YS}=\frac{3x}{4x}=\boxed{\mathrm{(B)}\ \frac{3}{4}}$



	\item (\textit{AMC 12 12A 2007}) A \href{https://artofproblemsolving.com/wiki/index.php?title=Triangle}{triangle} with side lengths in the \href{https://artofproblemsolving.com/wiki/index.php?title=Ratio}{ratio} $3 : 4 : 5$ is inscribed in a \href{https://artofproblemsolving.com/wiki/index.php?title=Circle}{circle} with \href{https://artofproblemsolving.com/wiki/index.php?title=Radius}{radius} 3. What is the area of the triangle?


 $\mathrm{(A)}\ 8.64\qquad \mathrm{(B)}\ 12\qquad \mathrm{(C)}\ 5\pi\qquad \mathrm{(D)}\ 17.28\qquad \mathrm{(E)}\ 18$

 \textbf{Answer: }A



\textbf{Solution}


                    \begin{figure}[H]
                        \centering
                        \includegraphics[width=0.5\textwidth]{Solution/AMC_12_12A_2007_10_Solution_Figure_1.png}
                    \end{figure}
                    

 Since 3-4-5 is a \href{https://artofproblemsolving.com/wiki/index.php?title=Pythagorean\_triple}{Pythagorean triple}, the triangle is a \href{https://artofproblemsolving.com/wiki/index.php?title=Right\_triangle}{right triangle}. Since the hypotenuse is a \href{https://artofproblemsolving.com/wiki/index.php?title=Diameter}{diameter} of the \href{https://artofproblemsolving.com/wiki/index.php?title=Circumcircle}{circumcircle}, the hypotenuse is $2r = 6$. Then the other legs are $\frac{24}5=4.8$ and $\frac{18}5=3.6$. The area is $\frac{4.8 \cdot 3.6}2 = 8.64\ \mathrm{(A)}$




\textbf{Solution 2}

The hypotenuse of the triangle is a diameter of the circumcircle, so it has length $2 \cdot 3 = 6$. The triangle is similar to a 3-4-5 triangle with the ratio of their side lengths equal to $\frac{6}{5}$. The area of a 3-4-5 triangle is $\frac{3\cdot 4}{2} = 6$.The square of the ratio of their side lengths is equal to the ratio of their areas. Call the area of the triangle $A$. Therefore, $\left(\frac{6}{5}\right)^2 = \frac{A}{6} \Longrightarrow \frac{36}{25} = \frac{A}{6} \Longrightarrow A = \frac{36\cdot6}{25} = 8.64\ \mathrm{(A)}$




\textbf{Solution 3 (dimensional analysis)}

Using the theorem that the hypotenuse of a right triangle inscribed in a circle is the diameter, one can determine that the diameter of the triangle is 6. As a result the ratio between the triangular ratios and their real values is (6/5). Using dimensional analysis, one can see that because the ratio between lengths is (6/5), the ratio between areas is (6/5)\textasciicircum{}2. This gives ((3 * 4)(6/5)\textasciicircum{}2)/2 as the area of the triangle or answer choice (A).

 



\textbf{Solution 4}

Let the sides of the triangle be $3x, 4x, 5x$ for some $x$. Now use the fact that the area of a triangle is equal to $\frac{abc}{4R}$, where $a, b, c$ are side lengths and $R$ is the circumradius. Also note that due to the 3, 4, 5 triple, the triangle is right. Thus we get $\frac{12x^2}{2} = \frac{60x^3}{4\cdot 3}\Rightarrow x = \frac{6}{5} = 1.2$. Then the area is $12(1.2)^2/2 = \boxed{8.64 \text{ (A)}}$.

 


	\item (\textit{AMC 12 12A 2007}) A finite \href{https://artofproblemsolving.com/wiki/index.php?title=Sequence}{sequence} of three-digit integers has the property that the tens and units digits of each term are, respectively, the hundreds and tens digits of the next term, and the tens and units digits of the last term are, respectively, the hundreds and tens digits of the first term. For example, such a sequence might begin with the terms 247, 475, and 756 and end with the term 824. Let $S$ be the sum of all the terms in the sequence. What is the largest \href{https://artofproblemsolving.com/wiki/index.php?title=Prime}{prime} \href{https://artofproblemsolving.com/wiki/index.php?title=Factor}{factor} that always divides $S$? 


 $\mathrm{(A)}\ 3\qquad \mathrm{(B)}\ 7\qquad \mathrm{(C)}\ 13\qquad \mathrm{(D)}\ 37\qquad \mathrm{(E)}\ 43$

 \textbf{Answer: }D



\textbf{Solution}

A given digit appears as the hundreds digit, the tens digit, and the units digit of a term the same number of times. Let $k$ be the sum of the units digits in all the terms. Then $S=111k=3 \cdot 37k$, so $S$ must be divisible by $37\ \mathrm{(D)}$. To see that it need not be divisible by any larger prime, the sequence $123, 231, 312$ gives $S=666=2 \cdot 3^2 \cdot 37\Rightarrow \mathrm{\boxed{(D)~~37}}$.

 



\textbf{Solution 2 (Principle of Fractions)}

Let the first term be $100a+10b+c$. The next term, $100b+10c+d$, then $100c+10d+e$, then $100d+10e+f$, and so on until $100w+10x+y$ and then $100x+10y+z$. Then, the sum $S$ is just the first term $100a$, added with $110b$, and then added with $111(c+d+\dots+x)$, followed by $11y$ and finally $z$. In mathematical terms,
 \[S = 100a+110b+111(c+d+\dots+x) + 11y + z\].

 The largest prime $p~|~S$ satisfies the property that it must equivalate itself as an integer when divided and added to all other terms or must divide one of the coefficients. If the former is true, then $100+110+111+11+1 = 333$ must equivalate to zero modulo that prime, implying that 3 and 37 are possible equivalencies. To verify no other cases, we check the latter case. Out of all plausible prime divisors of each respective coefficient we see that the extremum $111$ contains the largest, which is 37. We check modulo 37 on $S$ but quickly arise the same conclusion as the prior case, supporting that 37 is indeed the largest prime.

 The requested largest prime is $\boxed{37}$. $\square$

 \textbf{Post Problem Note} Note that here, $y$ and $z$ are interchangeable for $a$ and $b$ given the problem directory, and by substitution one will arrive to the same conclusion as solution 1.

 



\textbf{Solution 3 (Using the provided sequence/Guessing)}

Since a prime number will always divide $S$, the prime number must work on any sequence that follows the stated rule. Since the solution kindly provides us with the terms $247,475,756,824$,all we have to do is finish the sequence.

 One way of finishing the sequence is by adding $568,682$ (before 824) to the sequence and adding them all together. After this, we just divide the result by each of the answer choices. Of the five, $37$ is the largest prime number that divides this sequence, so the answer is ${\boxed{(D)}}$.~Banspeedrun

 



\textbf{Solution 4 (a better way of doing solution 3)}

Because the question says which prime number will \textbf{ \textit{ ALWAYS }} divide $S$, the prime number would divide any sequence that follows the rule. So, to test this, create two sequences with three terms each- $111, 111, 111$ and $222, 222, 222$. The sum for both (respectively) is $333$ and $666$. Obviously, $3$ divides both. Now, test the other choices on both numbers, but don't if it doesn't divide one. An example is $\frac{333}{7}$. Because it doesn't divide, don't check $\frac{666}{7}$. Use the same logic for $\frac{333}{13}$. However, when you test $\frac{333}{37}$, it works. Now, you must check $\frac{666}{37}$ to ensure it isn't a special case. After doing this, you get that both are divisible. Now, check $\frac{333}{43}$ and see that it doesn't work. 

 
So the answer is ${\boxed{(D)}}$

 -N.

 


	\item (\textit{AMC 12 12A 2007}) Integers $a, b, c,$ and $d$, not necessarily distinct, are chosen independently and at random from 0 to 2007, inclusive. What is the \href{https://artofproblemsolving.com/wiki/index.php?title=Probability}{probability} that $ad-bc$ is \href{https://artofproblemsolving.com/wiki/index.php?title=Even}{even}?


 $\mathrm{(A)}\ \frac 38\qquad \mathrm{(B)}\ \frac 7{16}\qquad \mathrm{(C)}\ \frac 12\qquad \mathrm{(D)}\ \frac 9{16}\qquad \mathrm{(E)}\ \frac 58$

 \textbf{Answer: }E



\textbf{Solution 1}

The only time when $ad-bc$ is even is when $ad$ and $bc$ are of the same \href{https://artofproblemsolving.com/wiki/index.php?title=Parity}{parity}. The chance of $ad$ being odd is $\frac 12 \cdot \frac 12 = \frac 14$, since the only way to have $ad$ be odd is to have both $a$ and $d$ be odd. As a result, $ad$ has a $\frac 34$ probability of being even. $bc$ also has a $\frac 14$ chance of being odd and a $\frac34$ chance of being even. Therefore, the probability that $ad-bc$ will be even is $\left(\frac 14\right)^2+\left(\frac 34\right)^2=\boxed {\mathrm{(E )} \frac 58\ }$.

 



\textbf{Solution 2 (casework)}

If we don't know our parity rules, we can check and see that $ad-bc$ is only even when $ad$ and $bc$ are of the same \href{https://artofproblemsolving.com/wiki/index.php?title=Parity}{parity} (as stated above). From here, we have two cases.

 Case 1: $odd-odd$ (which must be $o \cdot o-o \cdot o$). The probability for this to occur is $\left(\frac 12\right)^4 = \frac 1{16}$, because each integer has a $\frac 12$ chance of being odd.

 Case 2: $even-even$ (which occurs in 4 cases: $(e \cdot e-e \cdot e$), ($o \cdot e-o \cdot e$) (alternating of any kind), and ($e \cdot e-o \cdot e$) with its reverse, ($o \cdot e-e \cdot e$).

 Our first subcase of case 2 has a chance of $\frac 1{16}$ (same reasoning as above).

 Our second subcase of case 2 has a $\frac 14$ chance, since only the 2nd and 4th flip matter (or 1st and 3rd).

 Our third subcase of case 2 has a $\frac 18$ chance, because the 1st, 2nd, and either 3rd or 4th flip matter.

 Our fourth subcase of case 2 has a $\frac 18$ chance, because it's the same, just reversed.

 We sum these, and get our answer of $\frac 1{16} + \frac 1{16} + \frac 14 + \frac 18 + \frac 18 = \boxed {\mathrm{(E )} \frac 58\ }.$



	\item (\textit{AMC 12 12A 2007}) A piece of cheese is located at $(12,10)$ in a \href{https://artofproblemsolving.com/wiki/index.php?title=Coordinate\_plane}{coordinate plane}. A mouse is at $(4,-2)$ and is running up the \href{https://artofproblemsolving.com/wiki/index.php?title=Line}{line} $y=-5x+18$. At the point $(a,b)$ the mouse starts getting farther from the cheese rather than closer to it. What is  $a+b$?


 $\mathrm{(A)}\ 6\qquad \mathrm{(B)}\ 10\qquad \mathrm{(C)}\ 14\qquad \mathrm{(D)}\ 18\qquad \mathrm{(E)}\ 22$

 \textbf{Answer: }B



\textbf{Solution 1}

The point $(a,b)$ is the foot of the perpendicular from $(12,10)$ to the line $y=-5x+18$. The perpendicular has slope $\frac{1}{5}$, so its equation is $y=10+\frac{1}{5}(x-12)=\frac{1}{5}x+\frac{38}{5}$. The $x$-coordinate at the foot of the perpendicular satisfies the equation $\frac{1}{5}x+\frac{38}{5}=-5x+18$, so $x=2$ and $y=-5\cdot2+18=8$. Thus $(a,b) = (2,8)$, and $a+b = \boxed{10}$.

 



\textbf{Solution 2}

If the mouse is at $(x, y) = (x, 18 - 5x)$, then the square of the distance from the mouse to the cheese is $(x - 12)^2 + (8 - 5x)^2 = 26(x^2 - 4x + 8) = 26((x - 2)^2 + 4).$The value of this expression is smallest when $x = 2$, so the mouse is closest to the cheese at the point $(2, 8)$, and $a+b=2+8 = \boxed{10}$.

 



\textbf{Solution 3}

We are trying to find the point where distance between the mouse and $(12, 10)$ is minimized. This point is where the line that passes through $(12, 10)$ and is perpendicular to $y=-5x+18$ intersects $y=-5x+18$. By basic knowledge of perpendicular lines, this line is $y=\frac{x}{5}+\frac{38}{5}$. This line intersects $y=-5x+18$ at $(2,8)$. So $a+b=\boxed{10}$. - MegaLucario1001

 


 



\textbf{Solution 4}

If the mouse is at $(x, y) = (x, 18 - 5x)$, then the square of the distance from the mouse to the cheese is$(x - 12)^2 + (8 - 5x)^2 = 26(x^2 - 4x + 8) = 26((x - 2)^2 + 4)$.The value of this expression is smallest when $x = 2$, so the mouse is closest to the cheese at the point $(2, 8)$, and $a+b=2+8 = \boxed{10}$.-Paixiao

 


	\item (\textit{AMC 12 12A 2007}) Let a, b, c, d, and e be distinct \href{https://artofproblemsolving.com/wiki/index.php?title=Integer}{integers} such that


 $(6-a)(6-b)(6-c)(6-d)(6-e)=45$

 What is $a+b+c+d+e$?


 $\mathrm{(A)}\ 5\qquad \mathrm{(B)}\ 17\qquad \mathrm{(C)}\ 25\qquad \mathrm{(D)}\ 27\qquad \mathrm{(E)}\ 30$

 \textbf{Answer: }C



\textbf{Solution 1}

If $45$ is expressed as a product of five distinct integer factors, the absolute value of the product of any four is at least $|(-3)(-1)(1)(3)|=9$, so no factor can have an absolute value greater than $5$. Thus the factors of the given expression are five of the integers $\pm 3, \pm 1, \pm 5$. The product of all six of these is $-225=(-5)(45)$, so the factors are $-3, -1, 1, 3,$ and $5.$ The corresponding values of $a, b, c, d,$ and $e$ are $9, 7, 5, 3,$ and $1,$ and their sum is $\fbox{25 (C)}$




\textbf{Solution 2}

The prime factorization of $45$ is $3^2 * 5$. Therefore, the 5 distinct integer factors must have some negative numbers in them. Because there are two $3$'s in the prime factorization, one of them must be negative and the other positive. Because there is a $-3$, there must also be a $-1$ to cancel the negatives out. The 5 distinct integer factors must be $-3, 3, 5, -1, 1$. The corresponding values of $a, b, c, d,$ and $e$ are $9, 3, 1, 7, 5$. and their sum is $\fbox{25 (C)}$



	\item (\textit{AMC 12 12A 2007}) The \href{https://artofproblemsolving.com/wiki/index.php?title=Set}{set} $\{3,6,9,10\}$ is augmented by a fifth element $n$, not equal to any of the other four. The \href{https://artofproblemsolving.com/wiki/index.php?title=Median}{median} of the resulting set is equal to its \href{https://artofproblemsolving.com/wiki/index.php?title=Mean}{mean}. What is the sum of all possible values of $n$?


 $\mathrm{(A)}\ 7\qquad \mathrm{(B)}\ 9\qquad \mathrm{(C)}\ 19\qquad \mathrm{(D)}\ 24\qquad \mathrm{(E)}\ 26$

 \textbf{Answer: }E



\textbf{Solution}

The median must either be $6, 9,$ or $n$. \href{https://artofproblemsolving.com/wiki/index.php?title=Casework}{Casework}:

 
            \begin{itemize}
        \item Median is $6$: Then $n \le 6$ and $\frac{3+6+9+10+n}{5} = 6 \Longrightarrow n = 2$.
            \end{itemize}
            
            \begin{itemize}
        \item Median is $9$: Then $n \ge 9$ and $\frac{3+6+9+10+n}{5} = 9 \Longrightarrow n = 17$.
            \end{itemize}
            
            \begin{itemize}
        \item Median is $n$: Then $6 < n < 9$ and $\frac{3+6+9+10+n}{5} = n \Longrightarrow n = 7$.
            \end{itemize}
            All three cases are valid, so our solution is $2 + 7 + 17 = 26 \Longrightarrow \mathrm{(E)}$. 

 


	\item (\textit{AMC 12 12A 2007}) How many three-digit numbers are composed of three distinct digits such that one digit is the \href{https://artofproblemsolving.com/wiki/index.php?title=Average}{average} of the other two? 


 $\mathrm{(A)}\ 96\qquad \mathrm{(B)}\ 104\qquad \mathrm{(C)}\ 112\qquad \mathrm{(D)}\ 120\qquad \mathrm{(E)}\ 256$

 \textbf{Answer: }C



\textbf{Solution 1}

We can find the number of increasing \href{https://artofproblemsolving.com/wiki/index.php?title=Arithmetic\_sequence}{arithmetic sequences} of length 3 possible from 0 to 9, and then find all the possible permutations of these sequences.

 


	\item (\textit{AMC 12 12A 2007}) Suppose that $\sin a + \sin b = \sqrt{\frac{5}{3}}$ and $\cos a + \cos b = 1$. What is $\cos (a - b)$?


 $\mathrm{(A)}\ \sqrt{\frac{5}{3}} - 1\qquad \mathrm{(B)}\ \frac 13\qquad \mathrm{(C)}\ \frac 12\qquad \mathrm{(D)}\ \frac 23\qquad \mathrm{(E)}\ 1$

 \textbf{Answer: }B



\textbf{Solution}

We can make use the of the trigonometric \href{https://artofproblemsolving.com/wiki/index.php?title=Trigonometric\_identities\#Pythagorean\_Identities}{Pythagorean identities}: square both equations and add them up:

 $\sin^2 a + \sin^2 b + 2\sin a \sin b + \cos^2 a + \cos^2 b + 2\cos a \cos b = \frac{5}{3} + 1$
$2 + 2\sin a \sin b + 2\cos a \cos b = \frac{8}{3}$
$2(\cos a \cos b + \sin a \sin b) = \frac{2}{3}$This is just the cosine difference identity, which simplifies to $\cos (a - b) = \frac{1}{3} \Longrightarrow \mathrm{(B)}$



	\item (\textit{AMC 12 12A 2007}) The \href{https://artofproblemsolving.com/wiki/index.php?title=Polynomial}{polynomial} $f(x) = x^{4} + ax^{3} + bx^{2} + cx + d$ has real \href{https://artofproblemsolving.com/wiki/index.php?title=Coefficient}{coefficients}, and $f(2i) = f(2 + i) = 0.$ What is $a + b + c + d?$
$\mathrm{(A)}\ 0 \qquad \mathrm{(B)}\ 1 \qquad \mathrm{(C)}\ 4 \qquad \mathrm{(D)}\ 9 \qquad \mathrm{(E)}\ 16$

 \textbf{Answer: }D



\textbf{Solution 1}

A fourth degree polynomial has four \href{https://artofproblemsolving.com/wiki/index.php?title=Root}{roots}. Since the coefficients are real(meaning that complex roots come in conjugate pairs), the remaining two roots must be the \href{https://artofproblemsolving.com/wiki/index.php?title=Complex\_conjugate}{complex conjugates} of the two given roots. By the factor theorem, our roots are $2-i,-2i$. Now we work backwards for the polynomial:

 <div style="text-align:center;">$(x-(2+i))(x-(2-i))(x-2i)(x+2i) = 0$
$(x^2 - 4x + 5)(x^2 + 4) = 0$
$x^4 - 4x^3 + 9x^2 - 16x + 20 = 0$</div>Thus our answer is $- 4 + 9 - 16 + 20 = 9\ \mathrm{(D)}$.

 



\textbf{Solution 2}

Just like in Solution 1 we realize that the roots come in conjugate pairs. Which means the roots are $2i,i+2,-2i,2-i$So our polynomial is  

 \textbf{(1)} $f(x) = (x-2i)(x+2i)(x-i-2)(x-2+i)$

 Looking at the equation of the polynomial $f(x) = x^4+ax^3+bx^2+cx+d$. We see that $a+b+c+d = f(1)-1$

 If we plug in $1$ into equation \textbf{(1)} we get $f(1) = (1-2i)(1+2i)(-1-i)(-1+i)$. 

 Now if we multiply a complex number by its conjugate we get the sum of the squares of its real and imaginary parts. Using this property on the above we multiply  and get $f(1) = (1-2i)(1+2i)(-1-i)(-1+i) = (1^2+2^2)(1^2+1^2) = 10$So the answer is $f(1) -1 = 10 - 1  = 9$. $\framebox{D}$



	\item (\textit{AMC 12 12A 2007}) \href{https://artofproblemsolving.com/wiki/index.php?title=Triangle}{Triangles} $ABC$ and $ADE$ have \href{https://artofproblemsolving.com/wiki/index.php?title=Area}{areas} $2007$ and $7002,$ respectively, with $B = (0,0),$ $C = (223,0),$ $D = (680,380),$ and $E = (689,389).$ What is the sum of all possible x-coordinates of $A$?


 $\mathrm{(A)}\ 282 \qquad \mathrm{(B)}\ 300 \qquad \mathrm{(C)}\ 600 \qquad \mathrm{(D)}\ 900 \qquad \mathrm{(E)}\ 1200$

 \textbf{Answer: }E



\textbf{Solution}


                    \begin{figure}[H]
                        \centering
                        \includegraphics[width=0.5\textwidth]{Solution/AMC_12_12A_2007_19_Solution_Figure_1.png}
                    \end{figure}
                    




\textbf{Solution 1}

From $k = [ABC] = \frac 12bh$, we have that the height of $\triangle ABC$ is $h = \frac{2k}{b} = \frac{2007 \cdot 2}{223} = 18$. Thus $A$ lies on the lines $y = \pm 18 \quad \mathrm{(1)}$.

 $DE = 9\sqrt{2}$ using 45-45-90 triangles, so in $\triangle ADE$ we have that $h = \frac{2 \cdot 7002}{9\sqrt{2}} = 778\sqrt{2}$. The slope of $DE$ is $1$, so the equation of the line is $y = x + b \Longrightarrow b = (380) - (680) = -300 \Longrightarrow y = x - 300$. The point $A$ lies on one of two \href{https://artofproblemsolving.com/wiki/index.php?title=Parallel}{parallel} lines that are $778\sqrt{2}$ units away from $\overline{DE}$. Now take an arbitrary point on the line $\overline{DE}$ and draw the \href{https://artofproblemsolving.com/wiki/index.php?title=Perpendicular}{perpendicular} to one of the parallel lines; then draw a line straight down from the same arbitrary point. These form a 45-45-90 $\triangle$, so the straight line down has a length of $778\sqrt{2} \cdot \sqrt{2} = 1556$. Now we note that the \href{https://artofproblemsolving.com/wiki/index.php?title=Y-intercept}{y-intercept} of the parallel lines is either $1556$ units above or below the y-intercept of line $\overline{DE}$; hence the equation of the parallel lines is $y = x - 300 \pm 1556 \Longrightarrow x = y + 300 \pm 1556 \quad \mathrm{(2)}$.

 We just need to find the intersections of these two lines and sum up the values of the x-coordinates. Substituting the $\mathrm{(1)}$ into $\mathrm{(2)}$, we get $x = \pm 18 + 300 \pm 1556 = 4(300) = 1200 \Longrightarrow \mathrm{(E)}$.

 



\textbf{Solution 2}

We are finding the intersection of two pairs of \href{https://artofproblemsolving.com/wiki/index.php?title=Parallel}{parallel} lines, which will form a \href{https://artofproblemsolving.com/wiki/index.php?title=Parallelogram}{parallelogram}. The \href{https://artofproblemsolving.com/wiki/index.php?title=Centroid}{centroid} of this parallelogram is just the intersection of $\overline{BC}$ and $\overline{DE}$, which can easily be calculated to be $(300,0)$. Now the sum of the x-coordinates is just $4(300) = 1200$.

 



\textbf{Solution 3 (Bashing but very straightforward)}

After we compute that the y-value can be either $y = \pm 18$ and realize there are four total values (each pair being equally spaced on their respective y-lines of $\pm 18$), we can use an easy application of the \href{https://artofproblemsolving.com/wiki/index.php?title=Shoelace\_Theorem}{Shoelace Theorem} to figure out the values of X. Since we already know the two distances (positive y and negative y) will be the same, then we can simply plug in y=18, compute the sum of the two corresponding x-values and multiply it by two to get our answer which is $2(600) = 1200 \Longrightarrow \mathrm{(E)}$

 - Zephyrica

 



\textbf{Solution 4 (intense bashing, similar to Solution 3)}

We can use the shoelace theorem to first find that the y-coordinate of $A$ can be $-18$ or $18$. Then we can apply shoelace again to find the $4$ possible x-coordinates, namely $-1274$, $-1238$, $1838$, and $1874$. Adding these up, we get $1200 \Longrightarrow \mathrm{(E)}$.

 


	\item (\textit{AMC 12 12A 2007}) Corners are sliced off a \href{https://artofproblemsolving.com/wiki/index.php?title=Unit\_cube}{unit cube} so that the six \href{https://artofproblemsolving.com/wiki/index.php?title=Face}{faces} each become regular \href{https://artofproblemsolving.com/wiki/index.php?title=Octagon}{octagons}. What is the total \href{https://artofproblemsolving.com/wiki/index.php?title=Volume}{volume} of the removed \href{https://artofproblemsolving.com/wiki/index.php?title=Tetrahedra}{tetrahedra}?


 $\mathrm{(A)}\ \frac{5\sqrt{2}-7}{3}\qquad \mathrm{(B)}\ \frac{10-7\sqrt{2}}{3}\qquad \mathrm{(C)}\ \frac{3-2\sqrt{2}}{3}\qquad \mathrm{(D)}\ \frac{8\sqrt{2}-11}{3}\qquad \mathrm{(E)}\ \frac{6-4\sqrt{2}}{3}$

 \textbf{Answer: }B



\textbf{Solution}


                    \begin{figure}[H]
                        \centering
                        \includegraphics[width=0.5\textwidth]{Solution/AMC_12_12A_2007_20_Solution_Figure_1.png}
                    \end{figure}
                    

 Since the sides of a \href{https://artofproblemsolving.com/wiki/index.php?title=Regular\_polygon}{regular polygon} are equal in length, we can call each side $x$. Examine one edge of the unit cube: each contains two slanted diagonal edges of an octagon and one straight edge. The diagonal edges form $45-45-90 \triangle$ \href{https://artofproblemsolving.com/wiki/index.php?title=Right\_triangle}{right triangles}, making the distance on the edge of the cube $\frac{x}{\sqrt{2}}$. Thus, $2 \cdot \frac{x}{\sqrt{2}} + x = 1$, and $x = \frac{1}{\sqrt{2} + 1} \cdot \left(\frac{\sqrt{2} - 1}{\sqrt{2} - 1}\right) = \sqrt{2} - 1$. 

 Each of the cut off corners is a \href{https://artofproblemsolving.com/wiki/index.php?title=Pyramid}{pyramid}, whose volume can be calculated by $V = \frac 13 Bh$. Use the base as one of the three \href{https://artofproblemsolving.com/wiki/index.php?title=Congruency\_(geometry)&action=edit&redlink=1}{congruent} \href{https://artofproblemsolving.com/wiki/index.php?title=Isosceles\_triangle}{isosceles triangles}, with the height being one of the edges of the pyramid that sits on the edges of the cube. The height is $\frac x{\sqrt{2}} = 1 - \frac{1}{\sqrt{2}}$. The base is a $45-45-90 \triangle$ with leg of length $1 - \frac{1}{\sqrt{2}}$, making its area $\frac 12\left(1 - \frac 1{\sqrt{2}}\right)^2 = \frac{3 - 2\sqrt{2}}4$. Plugging this in, we get that the volume of one of the tetrahedra is $\frac 13 \left(1 - \frac{1}{\sqrt{2}}\right)\left(\frac{3 - 2\sqrt{2}}4\right) = \frac{10 - 7\sqrt{2}}{24}$. Since there are 8 removed corners, we get an answer of $\frac{10 - 7\sqrt{2}}{3} \Longrightarrow \boxed{\mathrm{B}}$



	\item (\textit{AMC 12 12A 2007}) The sum of the \href{https://artofproblemsolving.com/wiki/index.php?title=Root}{zeros}, the product of the zeros, and the sum of the \href{https://artofproblemsolving.com/wiki/index.php?title=Coefficient}{coefficients} of the \href{https://artofproblemsolving.com/wiki/index.php?title=Function}{function} $f(x)=ax^{2}+bx+c$ are equal. Their common value must also be which of the following? 


 $\textrm{(A)}\ \textrm{the\ coefficient\ of\ }x^{2}~~~ \textrm{(B)}\ \textrm{the\ coefficient\ of\ }x$
$\textrm{(C)}\ \textrm{the\ y-intercept\ of\ the\ graph\ of\ }y=f(x)$
$\textrm{(D)}\ \textrm{one\ of\ the\ x-intercepts\ of\ the\ graph\ of\ }y=f(x)$
$\textrm{(E)}\ \textrm{the\ mean\ of\ the\ x-intercepts\ of\ the\ graph\ of\ }y=f(x)$

 \textbf{Answer: }A



\textbf{Solution}

By \href{https://artofproblemsolving.com/wiki/index.php?title=Vieta\%27s\_Formulas}{Vieta's Formulas}, the sum of the roots of a \href{https://artofproblemsolving.com/wiki/index.php?title=Quadratic\_equation}{quadratic equation} is $\frac {-b}a$,  the product of the zeros is $\frac ca$, and the sum of the coefficients is $a + b + c$. Setting equal the first two tells us that $\frac {-b}{a} = \frac ca \Rightarrow b = -c$. Thus, $a + b + c = a + b - b = a$, so the common value is also equal to the coefficient of $x^2 \Longrightarrow \fbox{\textrm{A}}$. 

 To disprove the others, note that:

 
            \begin{itemize}
        \item $\mathrm{B}$: then $b = \frac {-b}a$, which is not necessarily true.
        \item $\mathrm{C}$: the y-intercept is $c$, so $c = \frac ca$, not necessarily true.
        \item $\mathrm{D}$: an x-intercept of the graph is a root of the polynomial, but this excludes the other root.
        \item $\mathrm{E}$: the mean of the x-intercepts will be the sum of the roots of the quadratic divided by 2.
            \end{itemize}
            


	\item (\textit{AMC 12 12A 2007}) For each positive integer $n$, let $S(n)$ denote the sum of the digits of $n.$ For how many values of $n$ is $n + S(n) + S(S(n)) = 2007?$
$\mathrm{(A)}\ 1 \qquad \mathrm{(B)}\ 2 \qquad \mathrm{(C)}\ 3 \qquad \mathrm{(D)}\ 4 \qquad \mathrm{(E)}\ 5$

 \textbf{Answer: }D



\textbf{Solution 1}

For the sake of notation, let $T(n) = n + S(n) + S(S(n))$. Obviously $n<2007$. Then the maximum value of $S(n) + S(S(n))$ is when $n = 1999$, and the sum becomes $28 + 10 = 38$. So the minimum bound is $2007-38=1969$. We do \href{https://artofproblemsolving.com/wiki/index.php?title=Casework}{casework} upon the tens digit:

 Case 1: $196u \Longrightarrow u = 9$. Easy to directly disprove.

 Case 2: $197u$. $S(n) = 1 + 9 + 7 + u = 17 + u$, and $S(S(n)) = 8+u$ if $u \le 2$ and $S(S(n)) = 2 + (u-3) = u-1$ otherwise. 

 
            \begin{enumerate}[label=(\roman*)]
                \item Subcase a: $T = 1970 + u + 17 + u + 8 + u = 1995 + 3u = 2007 \Longrightarrow u = 4$. This exceeds our bounds, so no solution here.\item Subcase b: $T = 1970 + u + 17 + u + u - 1 = 1986 + 3u = 2007 \Longrightarrow u = 7$. First solution.
            \end{enumerate}
            Case 3: $198u$. $S(n) = 18 + u$, and $S(S(n)) = 9 + u$ if $u \le 1$ and $2 + (u-2) = u$ otherwise.

 
            \begin{enumerate}[label=(\roman*)]
                \item Subcase a: $T = 1980 + u + 18 + u + 9 + u = 2007 + 3u = 2007 \Longrightarrow u = 0$. Second solution.\item Subcase b: $T = 1980 + u + 18 + u + u = 1998 + 3u = 2007 \Longrightarrow u = 3$. Third solution.
            \end{enumerate}
            Case 4: $199u$. But $S(n) > 19$, and $n + S(n)$ clearly sum to $> 2007$.

 Case 5: $200u$. So $S(n) = 2 + u$ and $S(S(n)) = 2 + u$ (recall that $n < 2007$), and $2000 + u + 2 + u + 2 + u = 2004 + 3u = 2007 \Longrightarrow u = 1$. Fourth solution.

 In total we have $4 \mathrm{(D)}$ solutions, which are $1977, 1980, 1983,$ and $2001$.

 



\textbf{Solution 2}

Clearly, $n > 1950$. We can break this into three cases:

 Case 1: $n \geq 2000$

 
            \begin{enumerate}[label=(\roman*)]
                \item Inspection gives $n = 2001$.
            \end{enumerate}
            Case 2: $n < 2000$, $n = \overline{19xy}$, $x + y < 10$

 
            \begin{enumerate}[label=(\roman*)]
                \item If you set up an equation, it reduces to
            \end{enumerate}
            $4x + y = 32$

 
            \begin{enumerate}[label=(\roman*)]
                \item which has as its only solution satisfying the constraints $x = 8$, $y = 0$.
            \end{enumerate}
            Case 3: $n < 2000$, $n = \overline{19xy}$, $x + y \geq 10$

 
            \begin{enumerate}[label=(\roman*)]
                \item This reduces to
            \end{enumerate}
            
            \begin{enumerate}[label=(\roman*)]
                \item $4x + y = 35$. The only two solutions satisfying the constraints for this equation are $x = 7$, $y = 7$ and $x = 8$, $y = 3$.
            \end{enumerate}
            The solutions are thus $1977, 1980, 1983, 2001$ and the answer is $\mathrm{(D)}\  4$.

 



\textbf{Solution 3 (Fastest casework)}

It is well-known that $n \equiv S(n)\equiv S(S(n)) \pmod{9}.$ Substituting, we have that 
 \[n+n+n \equiv 2007 \pmod{9} \implies n \equiv 0 \pmod{3}.\] Since $n \leq 2007,$ we must have that $\max S(n)=1+9+9+9=28.$ Now, we list out the possible values for $S(n)$ in a table, noting that it is a multiple of $3$ because $n$ is a multiple of $3.$
$\begin{tabular}{c|c c c c c c c c c c}   S(n) & 0 & 3 & 6 & 9 & 12 & 15 & 18 & 21 & 24 & 27 \\ \end{tabular}$

 

 Then, we compute the corresponding values of $S(S(n)).$
$\begin{tabular}{c|c c c c c c c c c c}   S(n) & 0 & 3 & 6 & 9 & 12 & 15 & 18 & 21 & 24 & 27 \\ \hline S(S(n)) & 0 & 3 & 6 & 9 & 3 & 6 & 9 & 3 & 6 & 9 \\ \end{tabular}$

 

 Finally, we may compute the corresponding values of $n$ using the fact that $n=2007-S(n)-S(S(n)).$
$\begin{tabular}{c|c c c c c c c c c c}  n & 2007 & 2001 & 1995 & 1989 & 1992 & 1986 & 1980 & 1983 & 1977 & 1971 \\ \hline S(n) & 0 & 3 & 6 & 9 & 12 & 15 & 18 & 21 & 24 & 27 \\ \hline S(S(n)) & 0 & 3 & 6 & 9 & 3 & 6 & 9 & 3 & 6 & 9 \\ \end{tabular}$

 

 Notice how all conditions are designed to be satisfied except whether $S(n)$ is accurate with respect to $n.$ So, the only thing that remains is to check this. We may eliminate, for example, when $n=2007$ we have $S(n)=9$ while the table states that it is $0.$ Proceeding similarly, we obtain the following table.

 $\begin{tabular}{c|c c c c c c c c c c}  n & \cancel{2007} & \boxed{2001} & \cancel{1995} & \cancel{1989} & \cancel{1992} & \cancel{1986} & \boxed{1980} & \boxed{1983} & \boxed{1977} & \cancel{1971} \\ \hline  S(n) & 0 & 3 & 6 & 9 & 12 & 15 & 18 & 21 & 24 & 27 \\ \hline S(S(n)) & 0 & 3 & 6 & 9 & 3 & 6 & 9 & 3 & 6 & 9 \\ \end{tabular}$

 

 It follows that there are $\boxed{4}\implies \textbf{(D)}$ possible values for $n.$ ~samrocksnature

 



\textbf{Solution 4}

As in Solution 1, we note that $S(n)\leq 28$ and $S(S(n))\leq 10$. 
Obviously, $n\equiv S(n)\equiv S(S(n)) \pmod 9$.  
As $2007\equiv 0 \pmod 9$, this means that $n\bmod 9 \in\{0,3,6\}$, or equivalently that $n\equiv S(n)\equiv S(S(n))\equiv 0 \pmod 3$.

 Thus $S(S(n))\in\{3,6,9\}$. For each possible $S(S(n))$ we get three possible $S(n)$.  
(E. g., if $S(S(n))=6$, then $S(n)=x$ is a number such that $x\leq 28$ and $S(x)=6$, therefore $S(n)\in\{6,15,24\}$.)

 For each of these nine possibilities we compute $n_{?}$ as $2007-S(n)-S(S(n))$ and check whether $S(n_{?})=S(n)$.  
We'll find out that out of the 9 cases, in 4 the value $n_{?}$ has the correct sum of digits.  
This happens for $n_{?}\in \{ 1977, 1980, 1983, 2001 \}$.

 



\textbf{Solution 5}

Claim. The only positive integers $n$ that satisfy the condition are perfect multiples of $3$.

 Proof of claim:We examine the positive integers mod $9$. Here are the cases.

 Case 1. $n \equiv 1 \pmod 9$. Now, we examine $S(n)$ modulo $9$.Case 1.1. The tens digit of $n$ is different from the tens digit of the largest multiple of $9$ under $n$. (In other words, this means we will carry when adding from the perfect multiple of $9$ under $n$.)Observe that when we carry, i.e. Add $1$ onto $1989$ to obtain $1990$, the units digit decreases by $9$ while the tens digit increases by $1$. This means that the sum of the digits decreases by $8$ in total, and we have $-8 \equiv 1 \pmod 9$, so the "mod 9" of the sum increases by $1$. This means that, regardless of whether the sum carries or not, the modulo 9 of the sum of the digits always increases by $1$.

 Case 1.2. The tens digits are the same, which is trivial since the units digit just increases by $1$ which means that the sum is also equivalent to $1 \pmod 9$.

 This means that $S(n) \equiv 1 \pmod 9$ and similarly letting the next $n=S(n)$, $S(S(n)) \equiv 1 \pmod 9$. Summing these, we have $n+S(n)+S(S(n)) \equiv 3 \pmod 9$. Clearly, no integers of this form will satisfy the condition because $2007$ is a perfect multiple of $9$.

 Case 2. $n \equiv 2 \pmod 9$.

 In this case, we apply exactly the same argument. There is at most one carry, which means that the sum of the digits will always be congruent to $2$ mod $9$. Then we can apply similar arguments to get $S(n) \equiv 2 \pmod 9$ and $S(S(n)) \equiv 2 \pmod 9$, so adding gives $n+S(n)+S(S(n)) \equiv 6 \pmod 9$.

 It is trivial to see that for $n \equiv k \pmod 9$, for $0 \leq k \leq 8$, we must have $n+S(n)+S(S(n)) \equiv 3k \pmod 9$. Only when $k=0, 3, 6$ is $3k$ a multiple of $9$, which means that $n$ must be a multiple of $3$.

 Now, we find the integers. Again, consider two cases: Integers that are direct multiples of $9$ and integers that are multiples of $3$ but not $9$.

 Case 1. $n$ is a multiple of $9$. An integer of the form $\overline{20ab}$ will not work since the least such integer is $2007$ which already exceeds our bounds. Thus, we need only consider the integers of the form $\overline{19ab}$. The valid sums of the digits of $n$ are $18$ and $27$ in this case.

 Case 1.1. The sum of the digits is $18$. This means that $S(n)=18, S(S(n))=9$, so $n=2007-18-9=1980$. Clearly this number satisfies our constraints.

 Case 1.2. The sum of the digits is $27$. This means that $S(n)=27, S(S(n))=9$, ,so $n=2007-27-9=1971$. Since the sum of the digits of $1971$ is not $27$, this does not work.

 This means that there is $1$ integer in this case.

 Case 2. $n$ is a multiple of $3$, not $9$..Case 2.1. Integers of the form $\overline{20ab}$. Then $S(n)=3$ or $S(n)=6$; it is trivial to see that $S(n)=6$ exceeds our bounds, so $S(n)=3$ and $n=2007-6=2001$.

 Case 2.2. Integers of the form $\overline{19ab}$. Then $S(n)=12, 15, 21, 24$ and we consider each case separately.

 Case 2.2.1. Integers with $S(n)=12$. That means $n=2007-12-3=1992$ which clearly does not work.

 Case 2.2.2. Integers with $S(n)=15$. That means $n=2007-15-6=1986$ which also does not work

 Case 2.2.3. Integers with $S(n)=21$. That means $n=2007-21-3=1983$ which is valid.

 Case 2.2.4. Integers with $S(n)=24$. That means $n=2007-24-6=1977$ which is also valid.

 We have considered every case, so there are $\boxed{4}$ integers that satisfy the given condition.

 



\textbf{Solution 6 (Rigorous)}

Let the number of digits of $n$ be $m$. If $m = 5$, $n$ will already be greater than $2007$. Notice that $S(n)$ is always at most $9m$. Then if $m = 3$, $n$ will be at most $999$, $S(n)$ will be at most $27$, and $S(S(n))$ will be even smaller than $27$. Clearly we cannot reach a sum of $2007$, unless $m = 4$ (i.e. $n$ has $4$ digits).

 Then, let $n$ be a four digit number in the form $1000a + 100b + 10c + d$. Then $S(n) = a + b + c + d$.

 $S(S(n))$ is the sum of the digits of $a + b + c + d$. We can represent $S(S(n))$ as the sum of the tens digit and the ones digit of $S(n)$. The tens digit in the form of a decimal is

 
$\frac{a + b + c + d}{10}$.

 
To remove the decimal portion, we can simply take the floor of the expression,

 
$\lfloor\frac{a + b + c + d}{10}\rfloor$.

 
Now that we have expressed the tens digit, we can express the ones digit as $S(n) -10$ times the above expression, or

 
$a + b + c + d - 10\lfloor\frac{a + b + c + d}{10}\rfloor$.

 
Adding the two expressions yields the value of $S(S(n))$
$= a + b + c + d - 9\lfloor\frac{a + b + c + d}{10}\rfloor$.

 
Combining this expression to the ones for $n$ and $S(n)$ yields

 
$1002a + 102b + 12c + 3d - 9\lfloor\frac{a + b + c + d}{10}\rfloor$.

 
Setting this equal to $2007$ and rearranging a bit yields

 
$12c + 3d = 2007 - 1002a - 102b + 9\lfloor\frac{a + b + c + d}{10}\rfloor$
$\Rightarrow$ $4c + d = 669 - 334a - 34b + 3\lfloor\frac{a + b + c + d}{10}\rfloor$.

 
(The reason for this slightly weird arrangement will soon become evident)

 
Now we examine the possible values of $a$. If $a \ge 3$, $n$ is already too large. $a$ must also be greater than $0$, or $n$ would be a $3$-digit number. Therefore, $a = 1 \, \text{or} \, 2$. Now we examine by case.

 If $a = 2$, then $b$ and $c$ must both be $0$ (otherwise $n$ would already be greater than $2007$). Substituting these values into the equation yields

 
$d = 1 + 3\lfloor\frac{2 + d}{10}\rfloor$
$\Rightarrow$ $d=1$.

 
Sure enough, $2001 + (2+1) + 3=2007$.

 Now we move onto the case where $a = 1$. Then our initial equation simplifies to

 
$4c + d = 335 - 34b + 3\lfloor\frac{1 + b + c + d}{10}\rfloor$

 
Since $c$ and $d$ can each be at most $9$, we substitute that value to find the lower bound of $b$. Doing so yields

 
$34b \ge 290 + 3\lfloor\frac{19 + b}{10}\rfloor$.

 
The floor expression is at least $3\lfloor\frac{19}{10}\rfloor=3$ , so the right-hand side is at least $293$. Solving for $b$, we see that $b \ge 9$ $\Rightarrow$ $b=9$. Again, we substitute for $b$ and the equation becomes

 
$4c + d = 29 + 3\lfloor\frac{10 + c + d}{10}\rfloor$
$\Rightarrow$ $4c + d = 32 + 3\lfloor\frac{c + d}{10}\rfloor$.

 
Just like we did for $b$, we can find the lower bound of $c$ by assuming $d = 9$ and solving:

 
$4c + 9 \ge 29 + 3\lfloor\frac{c + 9}{10}\rfloor$
$\Rightarrow$ $4c \ge 20 + 3\lfloor\frac{c + 9}{10}\rfloor$

 
The right hand side is $20$ for $c=0$ and $23$ for $c \ge 1$. Solving for c yields $c \ge 6$. Looking back at the previous equation, the floor expression is $0$ for $c+d \le 9$ and $3$ for $c+d \ge 10$. Thus, the right-hand side is $32$ for $c+d \le 9$ and $35$ for $c+d \ge 10$. We can solve these two scenarios as systems of equations/inequalities:

 $4c+d = 32$
$c+d \le 9$

 and

 $4c+d=35$
$c+d \ge 10$

 Solving yields three pairs $(c, d):$ $(8, 0)$; $(8, 3)$; and $(7, 7)$. Checking the numbers $1980$, $1983$, and $1977$; we find that all three work. Therefore there are a total of $4$ possibilities for $n$ $\Rightarrow$ $\boxed{\text{D}}$.

 Note: Although this solution takes a while to read (as well as to write) the actual time it takes to think through the process above is very short in comparison to the solution length.

 


	\item (\textit{AMC 12 12A 2007}) \href{https://artofproblemsolving.com/wiki/index.php?title=Square}{Square} $ABCD$ has area $36,$ and $\overline{AB}$ is \href{https://artofproblemsolving.com/wiki/index.php?title=Parallel}{parallel} to the \href{https://artofproblemsolving.com/wiki/index.php?title=X-axis}{x-axis}. Vertices $A,$ $B$, and $C$ are on the graphs of $y = \log_{a}x,$ $y = 2\log_{a}x,$ and $y = 3\log_{a}x,$ respectively. What is $a?$
$\mathrm{(A)}\ \sqrt [6]{3}\qquad \mathrm{(B)}\ \sqrt {3}\qquad \mathrm{(C)}\ \sqrt [3]{6}\qquad \mathrm{(D)}\ \sqrt {6}\qquad \mathrm{(E)}\ 6$

 \textbf{Answer: }A



\textbf{Solution 1}

Let $x$ be the x-coordinate of $B$ and $C$, and $x_2$ be the x-coordinate of $A$ and $y$ be the y-coordinate of $A$ and $B$. Then $2\log_ax= y \Longrightarrow a^{y/2} = x$ and $\log_ax_2 = y \Longrightarrow x_2 = a^y = \left(a^{y/2}\right)^2 = x^2$. Since the distance between $A$ and $B$ is $6$, we have $x^2 - x - 6 = 0$, yielding $x = -2, 3$.

 However, we can discard the negative root (all three \href{https://artofproblemsolving.com/wiki/index.php?title=Logarithm}{logarithmic} equations are underneath the line $y = 3$ and above $y = 0$ when $x$ is negative, hence we can't squeeze in a square of side 6). Thus $x = 3$.

 Substituting back, $3\log_{a}x - 2\log_{a}x = 6 \Longrightarrow a^6 = x$, so $a = \sqrt[6]{3}\ \ \mathrm{(A)}$.

 



\textbf{Solution 2}

Notice that all of the graphs $y = \log_{a}x,$ $y = 2\log_{a}x,$ and $y = 3\log_{a}x$ have a domain of $(0, \infty]$. Also notice that $y = \log_{a}x$ is the furthest to the right, as adding coefficients in front of the $\log$ part only makes the graph steeper. Since $A$ is on the graph of $y = \log_{a}x$ and $B$ is on the graph of $y = 2\log_{a}x$, $\,$ $A$ must be to the right of $B$. We are told that $\overline{AB}$ is \href{https://artofproblemsolving.com/wiki/index.php?title=Parallel}{parallel} to the \href{https://artofproblemsolving.com/wiki/index.php?title=X-axis}{x-axis}. Let $A$ be the point $(x, y)$. Then the points $B$ and $C$ are

 $(x-6, y)$

 and

 $(x-6, y+6)$

 respectively.

 Substituting these coordinates into the equations given yields $y = \log_{a}x,$ $y = 2\log_{a}x-6,$ and $y+6 = 3\log_{a}x-6$. Rearranging a bit, we get the following equations:

 $1) a^y = x$
$2) a^y = (x-6)^2$
$3) a^{y+6} = (x-6)^3$

 Using equations 1 and 2, we get $x=(x-6)^2$. Solving yields $x=9, x=4$. However, $4$ is extraneous because $x-6=4-6=-2$ is negative and makes it impossible for points $B$ and $C$ to be on the graphs of $y = 2\log_{a}x$ and $y = 3\log_{a}x$, respectively. This means $x=9$ (also because we know $A$ has to have a positive $x$-coordinate). Using the second and third equations, we get

 $\frac{a^{y+6}}{a^y} = \frac{(x-6)^3}{(x-6)^2}$ $\Rightarrow$ $a^6 = x-6$.

 Plugging in $9$ for $x$ yields $a^6 = 3$, or $a= \sqrt[6]{3} \Rightarrow \boxed{\text{A}}$.

 


	\item (\textit{AMC 12 12A 2007}) For each integer $n>1$, let $F(n)$ be the number of solutions to the equation $\sin{x}=\sin{(nx)}$ on the interval $[0,\pi]$. What is $\sum_{n=2}^{2007} F(n)$?


 $\mathrm{(A)}\ 2014524 \qquad \mathrm{(B)}\ 2015028 \qquad \mathrm{(C)}\ 2015033 \qquad \mathrm{(D)}\ 2016532 \qquad \mathrm{(E)}\ 2017033$

 \textbf{Answer: }D



\textbf{Solution 1}

$F(2)=3$

 By looking at various graphs, we obtain that, for most of the graphs

 $F(n) = n + 1$

 Notice that the solutions are basically reflections across $x = \frac{\pi}{2}$.However, when $n \equiv 1 \pmod{4}$, the middle apex of the \href{https://artofproblemsolving.com/wiki/index.php?title=Sine}{sine} curve touches the sine curve at the top only one time (instead of two reflected points), so we get here $F(n) = n$.

 $3+4+5+5+7+8+9+9+\cdots+2008$$= (1+2+3+4+5+\cdots+2008) - 3 - 501$$= \frac{(2008)(2009)}{2} - 504 = 2016532$ $\mathrm{(D)}$




\textbf{Solution 2}


 \[\sin nx - \sin x = 2\left( \cos \frac {n + 1}{2}x\right) \left( \sin \frac {n - 1}{2}x\right)\]So $\sin nx = \sin x$ if and only if $\cos \frac {n + 1}{2}x = 0$ or $\sin \frac {n - 1}{2}x = 0$.

 The first occurs whenever $\frac {n + 1}{2}x = (j + 1/2)\pi$, or $x = \frac {(2j + 1)\pi}{n + 1}$ for some nonnegative integer $j$. Since $x\leq \pi$, $j\leq n/2$. So there are $1 + \lfloor n/2 \rfloor$ solutions in this case.

 The second occurs whenever $\frac {n - 1}{2}x = k\pi$, or $x = \frac {2k\pi}{n - 1}$ for some nonnegative integer $k$. Here $k\leq \frac {n - 1}{2}$ so that there are $\left\lfloor \frac {n + 1}{2}\right\rfloor$ solutions here.

 However, we overcount intersections. These occur whenever
 \[\frac {2j + 1}{n + 1} = \frac {2k}{n - 1}\]

 \[k = \frac {(2j + 1)(n - 1)}{2(n + 1)}\]which is equivalent to $2(n + 1)$ dividing $(2j + 1)(n - 1)$. If $n$ is even, then $(2j + 1)(n - 1)$ is odd, so this never happens. If $n\equiv 3\pmod{4}$, then there won't be intersections either, since a multiple of 8 can't divide a number which is not even a multiple of 4.

 This leaves $n\equiv 1\pmod{4}$. In this case, the divisibility becomes $\frac {n + 1}{2}$ dividing $(2j + 1)\frac {n - 1}{4}$. Since $\frac {n + 1}{2}$ and $\frac {n - 1}{4}$ are relatively prime (subtracting twice the second number from the first gives 1), $\frac {n + 1}{2}$ must divide $2j + 1$. Since $j\leq \frac {n - 1}{2}$, $2j + 1\leq n < 2\cdot \frac {n + 1}{2}$. Then there is only one intersection, namely when $j = \frac {n - 1}{4}$.

 Therefore we find $F(n)$ is equal to $1 + \lfloor n/2 \rfloor + \left \lfloor \frac {n + 1}{2}\right\rfloor = n + 1$, unless $n\equiv 1\pmod{4}$, in which case it is one less, or $n$. The problem may then be finished as in Solution 1.

 Note from Williamgolly:An easier way to see that there is an extra point of intersection at $n \equiv 1 \pmod{4}$ is that $v_2(2j+1)=0,$ so we must have $v_2(n-1)>v_2(n+1)$ since $v_2(2k) \geq 1 >v_2(2j+1).$  Therefore, we suspect that $n \equiv 1 \pmod{4}$ has an extra point of intersection.  Testing, we see this is true.

 


	\item (\textit{AMC 12 12A 2007}) Call a set of integers \textit{spacy} if it contains no more than one out of any three consecutive integers. How many \href{https://artofproblemsolving.com/wiki/index.php?title=Subset}{subsets} of $\{1,2,3,\ldots,12\},$ including the \href{https://artofproblemsolving.com/wiki/index.php?title=Empty\_set}{empty set}, are spacy?


 $\mathrm{(A)}\ 121 \qquad \mathrm{(B)}\ 123 \qquad \mathrm{(C)}\ 125 \qquad \mathrm{(D)}\ 127 \qquad \mathrm{(E)}\ 129$

 \textbf{Answer: }E



\textbf{Solution 1}

Let $S_{n}$ denote the number of spacy subsets of $\{ 1, 2, ... n \}$.  We have $S_{0} = 1, S_{1} = 2, S_{2} = 3$.  

 The spacy subsets of $S_{n + 1}$ can be divided into two groups:

 
            \begin{itemize}
        \item $A =$ those not containing $n + 1$.  Clearly $|A|=S_{n}$.
        \item $B =$ those containing $n + 1$. We have $|B|=S_{n - 2}$, since removing $n + 1$ from any set in $B$ produces a spacy set with all elements at most equal to $n - 2,$ and each such spacy set can be constructed from exactly one spacy set in $B$.
            \end{itemize}
            Hence,

 $S_{n + 1} = S_{n} + S_{n - 2}$From this \href{https://artofproblemsolving.com/wiki/index.php?title=Recursion}{recursion}, we find that

 


\end{enumerate}